\documentclass[12pt]{article}
 
\usepackage[margin=1in]{geometry} 
\usepackage{amsmath,amsfonts,amsthm,amssymb}
\usepackage{mathtools}
\usepackage{pifont}
\usepackage{graphics}
\usepackage{graphicx}
\usepackage{array}
\usepackage{listings}
\usepackage{matlab-prettifier}
\usepackage[hidelinks]{hyperref}
\usepackage{caption}
\usepackage{enumitem}

\lstdefinestyle{tex-style} {
	language=Matlab,                           % Code langugage
	basicstyle=\ttfamily,                   % Code font, Examples: \footnotesize, \ttfamily
	%keywordstyle=\color{OliveGreen},        % Keywords font ('*' = uppercase)
	commentstyle=\color{gray},              % Comments font
	numbers=none,                           % Line nums position
	numberstyle=\tiny,                      % Line-numbers fonts
	stepnumber=1,                           % Step between two line-numbers
	numbersep=5pt,                          % How far are line-numbers from code
	%backgroundcolor=\color{gray}, % Choose background color
	frame=single,                             % A frame around the code
	tabsize=2,                              % Default tab size
	captionpos=b,                           % Caption-position = bottom
	breaklines=true,                        % Automatic line breaking?
	breakatwhitespace=false,                % Automatic breaks only at whitespace?
	showspaces=false,                       % Dont make spaces visible
	showtabs=false,                         % Dont make tabls visible
	columns=flexible,                       % Column format
	morekeywords={__global__, __device__}  % CUDA specific keywords
}
\lstset{style=tex-style}
 
\newcommand{\N}{\mathbb{N}}
\newcommand{\Z}{\mathbb{Z}}
 
\newenvironment{theorem}[2][Theorem]{\begin{trivlist}
\item[\hskip \labelsep {\bfseries #1}\hskip \labelsep {\bfseries #2.}]}{\end{trivlist}}
\newenvironment{lemma}[2][Lemma]{\begin{trivlist}
\item[\hskip \labelsep {\bfseries #1}\hskip \labelsep {\bfseries #2.}]}{\end{trivlist}}
\newenvironment{exercise}[2][Exercise]{\begin{trivlist}
\item[\hskip \labelsep {\bfseries #1}\hskip \labelsep {\bfseries #2.}]}{\end{trivlist}}
\newenvironment{problem}[2][Problem]{\begin{trivlist}
\item[\hskip \labelsep {\bfseries #1}\hskip \labelsep {\bfseries #2.}]}{\end{trivlist}}
\newenvironment{question}[2][Question]{\begin{trivlist}
\item[\hskip \labelsep {\bfseries #1}\hskip \labelsep {\bfseries #2.}]}{\end{trivlist}}
\newenvironment{corollary}[2][Corollary]{\begin{trivlist}
\item[\hskip \labelsep {\bfseries #1}\hskip \labelsep {\bfseries #2.}]}{\end{trivlist}}

\newenvironment{solution}{\begin{proof}[Solution]}{\end{proof}}

%%%%%%% Table Packages %%%%%%%%%%%%%%%%%%%%%%%
\usepackage{booktabs}
\usepackage{multirow}
\usepackage{array}
\usepackage[T1]{fontenc}
\usepackage[utf8]{inputenc} 
\usepackage{tabularx,ragged2e,booktabs}
\newcolumntype{C}[1]{>{\Centering}m{#1}}
\renewcommand\tabularxcolumn[1]{C{#1}}
\newcommand{\ra}[1]{\renewcommand{\arraystretch}{#1}}
\usepackage{makecell}
\usepackage{tikz}
\usetikzlibrary{tikzmark,calc}
\usepackage{siunitx}

 
\begin{document}
 
% --------------------------------------------------------------
%                         Start here
% --------------------------------------------------------------
 
\title{Homework 3}
\author{MEGN544/ROBO554 Robot Mechanics: Kinematics, Dynamics, and Control  
\\ 
Dr. George Kontoudis\\
Department of Mechanical Engineering\\
Colorado School of Mines}
\date{Fall 2026}
 
\maketitle
\begin{problem}{1}%theorem, exercise, problem, or question 
Given the following twist $\boldsymbol{\xi}$, what is the homogeneous transformation matrix?
\begin{enumerate}[label=(\alph*)]
    \item $\boldsymbol{\xi} = \begin{bmatrix}
        \boldsymbol{v}^{\intercal} & \boldsymbol{\Omega}^{\intercal}
        \end{bmatrix}^{\intercal} = \begin{bmatrix}
        3/\pi & 0 & 0 & \pi/3 & 0 & 0
        \end{bmatrix}^{\intercal}$.
    \item $\boldsymbol{\xi} = \begin{bmatrix}
        \boldsymbol{v}^{\intercal} & \boldsymbol{\Omega}^{\intercal}
        \end{bmatrix}^{\intercal} = \begin{bmatrix}
        0 & 0.5 & 0 & 0 & \pi/2 & 0
        \end{bmatrix}^{\intercal}$.
    \item $\boldsymbol{\xi} = \begin{bmatrix}
        \boldsymbol{v}^{\intercal} & \boldsymbol{\Omega}^{\intercal}
        \end{bmatrix}^{\intercal} = \begin{bmatrix}
        1 & 0 & 0 & 0 & 0 & \pi/3
        \end{bmatrix}^{\intercal}$.
\end{enumerate}
 
\end{problem}

\begin{problem}{2}%theorem, exercise, problem, or question 
Given the homogeneous transformation matrix, 
\begin{align*}
        \boldsymbol{T} = \begin{bmatrix}
        0.5000 & -0.6124 & 0.6124 & 0.8415 \\
        0.6124 & 0.7500 & 0.2500 & 0.3251 \\
        -0.6124 & 0.2500 & 0.7500 & -0.3251\\
        0 & 0 & 0 & 1
        \end{bmatrix}
\end{align*}
\begin{enumerate}[label=(\alph*)]
    \item What is the displacement?
    \item What is the Quaternion that encodes the rotation?
    \item What is the Angle-Axis that encodes the rotation?
    \item What is the twist $\boldsymbol{\xi}$ that encodes the full transform?
\end{enumerate}
 
\end{problem}

\begin{problem}{3}%theorem, exercise, problem, or question 
For the following give one pro and one con for using it as a representation,
\begin{enumerate}[label=(\alph*)]
    \item Rotation and displacement separately.
    \item Homogeneous transformation matrices.
    \item Twist coordinates.
\end{enumerate}
 
\end{problem}

\begin{problem}{4}%theorem, exercise, problem, or question 
Given the following two homogeneous transformations, use MATLAB to plot a linear
interpolation between the two. Your plot should contain a minimum of $5$ intermediate points and should have the orthogonal axis (columns of the rotation matrix) drawn at each point. \textit{Hint}: Use the \verb|quiver3| function.

\begin{align*}\boldsymbol{T}_{\textrm{initial}} = \begin{bmatrix}
        1 & 0 & 0 & 0 \\
        0 & 1 & 0 & 0 \\
        0 & 0 & 1 & 0\\
        0 & 0 & 0 & 1
        \end{bmatrix},~\boldsymbol{T}_{\textrm{final}} = \begin{bmatrix}
        0.5000 & -0.6124 & 0.6124 & 0.8415 \\
        0.6124 & 0.7500 & 0.2500 & 0.3251 \\
        -0.6124 & 0.2500 & 0.7500 & -0.3251\\
        0 & 0 & 0 & 1
        \end{bmatrix}
\end{align*}

\begin{enumerate}[label=(\alph*)]
    \item Perform the interpolation using ZYZ angles and position.
    \item Perform the interpolation using Angle-Axis and position.
    \item Perform the interpolation using Quaternion and position (remember to enforce the length constraint on the quaternion before calculating the intermediate orientations).
    \item Perform the interpolation using twist.
    \item Comment on the similarities or differences.
\end{enumerate}

\end{problem}


\end{document}